% GENERATED FILE. Edit canonical lessons, then run npm run generate:books.
% canonical-source-hash: fnv1a64:e73534edb1da725a
% canonical-lessons: MW-C127-khilo, MW-C127-hathoro, MW-C127-kulhari, MW-C127-hal, MW-C127-bailgari

\chapter{Nail, Hammer, Axe, Plough, Bullock cart}
\label{ch:mw-nail-hammer-axe-plough-bullock-cart}
\hlchaptermodality{127}

\begin{chapteropening}
\textbf{By the end of this chapter:} \emph{I can say a nail, a hammer, an axe, a plough and a bullock cart.}

Complete the last lesson of chapter 127: I can say a nail, a hammer, an axe, a plough and a bullock cart.
\end{chapteropening}

\section[\texorpdfstring{\mw{खीलो} (khīlo)}{खीलो (khīlo)}]{\mw{खीलो} (khīlo) — a nail}

\label{lesson:MW-C127-khilo}



\begin{quote}
Before the new one: say the Marwadi for a rope, then the Marwadi for a peg.
\end{quote}

\subsection*{You'll want to know: \mw{खीलो}}
\textbf{\mw{खीलो}} — \emph{khīlo} — \textquotedblleft{}a nail\textquotedblright{}.

\begin{cousinweb}
\textbf{\mw{खीलो}} (\emph{khīlo}) is a tadbhava, a word inherited from Sanskrit and reshaped by sound change. It comes from \emph{kīla}, \textquotedblleft{}stake, peg\textquotedblright{}. CDIAL gives no Marwari form in this entry (CDIAL 3202); the match is through Gujarati \emph{khīlo}. CDIAL notes that the \emph{k}/\emph{kh} variation points to a non-Aryan source.
\end{cousinweb}

\begin{grammarlens}[title={What you've built}]
One more word for this chapter.
\end{grammarlens}

\subsection*{Your turn}
\begin{itemize}
\raggedright
  \item \emph{Say it:} \emph{khīlo}
  \item \emph{Say it:} \emph{khīlo}, once more
  \item \emph{Say it:} say \emph{khīlo}
  \item \emph{Recall:} say the Marwadi for a rope, then the Marwadi for a peg, then say \emph{khīlo} again
\end{itemize}

\begin{tcolorbox}[breakable,colback=teal!4,colframe=teal!35!black,title={Before you move on}]
What does \mw{खीलो} mean? (\textquotedblleft{}A nail\textquotedblright{}.) Say it once more.
\end{tcolorbox}
\section[\texorpdfstring{\mw{हथोड़ो} (hathoṛo)}{हथोड़ो (hathoṛo)}]{\mw{हथोड़ो} (hathoṛo) — a hammer}

\label{lesson:MW-C127-hathoro}



\begin{quote}
Before the new one: say the Marwadi for a peg, then the Marwadi for a nail.
\end{quote}

\subsection*{You'll want to know: \mw{हथोड़ो}}
\textbf{\mw{हथोड़ो}} — \emph{hathoṛo} — \textquotedblleft{}a hammer\textquotedblright{}.

\begin{cousinweb}
\textbf{\mw{हथोड़ो}} (\emph{hathoṛo}) is a tadbhava, a word inherited and reshaped by sound change. It comes from \emph{*hastakūṭa}, \textquotedblleft{}hand hammer\textquotedblright{}, a form CDIAL reconstructs from its descendants. CDIAL gives no Marwari form in this entry (CDIAL 14028); the match is through Hindi \emph{hathauṛā} and Gujarati \emph{hathoṛo}. It is built on \emph{hasta}, \textquotedblleft{}hand\textquotedblright{}, and \emph{kūṭa}, \textquotedblleft{}hammer\textquotedblright{}.
\end{cousinweb}

\begin{grammarlens}[title={What you've built}]
One more word for this chapter.
\end{grammarlens}

\subsection*{Your turn}
\begin{itemize}
\raggedright
  \item \emph{Say it:} \emph{hathoṛo}
  \item \emph{Say it:} \emph{hathoṛo}, once more
  \item \emph{Say it:} say \emph{hathoṛo}
  \item \emph{Recall:} say the Marwadi for a peg, then the Marwadi for a nail, then say \emph{hathoṛo} again
\end{itemize}

\begin{tcolorbox}[breakable,colback=teal!4,colframe=teal!35!black,title={Before you move on}]
What does \mw{हथोड़ो} mean? (\textquotedblleft{}A hammer\textquotedblright{}.) Say it once more.
\end{tcolorbox}
\section[\texorpdfstring{\mw{कुल्हाड़ी} (kulhāṛī)}{कुल्हाड़ी (kulhāṛī)}]{\mw{कुल्हाड़ी} (kulhāṛī) — an axe}

\label{lesson:MW-C127-kulhari}



\begin{quote}
Before the new one: say the Marwadi for a nail, then the Marwadi for a hammer.
\end{quote}

\subsection*{You'll want to know: \mw{कुल्हाड़ी}}
\textbf{\mw{कुल्हाड़ी}} — \emph{kulhāṛī} — \textquotedblleft{}an axe\textquotedblright{}.

\begin{cousinweb}
\textbf{\mw{कुल्हाड़ी}} (\emph{kulhāṛī}) is a tadbhava, a word inherited from Sanskrit and reshaped by sound change. It comes from \emph{kuṭhāra}, \textquotedblleft{}axe\textquotedblright{}. CDIAL gives no Marwari form in this entry (CDIAL 3244); the match is through Hindi \emph{kulhāṛī}. CDIAL thinks the Sanskrit word itself probably came from Dravidian.
\end{cousinweb}

\begin{grammarlens}[title={What you've built}]
One more word for this chapter.
\end{grammarlens}

\subsection*{Your turn}
\begin{itemize}
\raggedright
  \item \emph{Say it:} \emph{kulhāṛī}
  \item \emph{Say it:} \emph{kulhāṛī}, once more
  \item \emph{Say it:} say \emph{kulhāṛī}
  \item \emph{Recall:} say the Marwadi for a nail, then the Marwadi for a hammer, then say \emph{kulhāṛī} again
\end{itemize}

\begin{tcolorbox}[breakable,colback=teal!4,colframe=teal!35!black,title={Before you move on}]
What does \mw{कुल्हाड़ी} mean? (\textquotedblleft{}An axe\textquotedblright{}.) Say it once more.
\end{tcolorbox}
\section[\texorpdfstring{\mw{हळ} (haḷ)}{हळ (haḷ)}]{\mw{हळ} (haḷ) — a plough}

\label{lesson:MW-C127-hal}



\begin{quote}
Before the new one: say the Marwadi for a hammer, then the Marwadi for an axe.
\end{quote}

\subsection*{You'll want to know: \mw{हळ}}
\textbf{\mw{हळ}} — \emph{haḷ} — \textquotedblleft{}a plough\textquotedblright{}.

\begin{cousinweb}
\textbf{\mw{हळ}} (\emph{haḷ}) is a tadbhava, a word inherited from Sanskrit and reshaped by sound change. It comes from \emph{hala}, \textquotedblleft{}plough\textquotedblright{}. CDIAL gives no Marwari form in this entry (CDIAL 14000); the match is through Hindi \emph{hal} and Gujarati \emph{haḷ}. Marwadi, like Gujarati, says the \emph{l} curled back.
\end{cousinweb}

\begin{grammarlens}[title={What you've built}]
One more word for this chapter.
\end{grammarlens}

\subsection*{Your turn}
\begin{itemize}
\raggedright
  \item \emph{Say it:} \emph{haḷ}
  \item \emph{Say it:} \emph{haḷ}, once more
  \item \emph{Say it:} say \emph{haḷ}
  \item \emph{Recall:} say the Marwadi for a hammer, then the Marwadi for an axe, then say \emph{haḷ} again
\end{itemize}

\begin{tcolorbox}[breakable,colback=teal!4,colframe=teal!35!black,title={Before you move on}]
What does \mw{हळ} mean? (\textquotedblleft{}A plough\textquotedblright{}.) Say it once more.
\end{tcolorbox}
\section[\texorpdfstring{\mw{बैलगाड़ी} (bailgāṛī)}{बैलगाड़ी (bailgāṛī)}]{\mw{बैलगाड़ी} (bailgāṛī) — a bullock cart}

\label{lesson:MW-C127-bailgari}



\begin{quote}
Before the new one: say the Marwadi for an axe, then the Marwadi for a plough.
\end{quote}

\subsection*{You'll want to know: \mw{बैलगाड़ी}}
\textbf{\mw{बैलगाड़ी}} — \emph{bailgāṛī} — \textquotedblleft{}a bullock cart\textquotedblright{}.

\begin{cousinweb}
\textbf{\mw{बैलगाड़ी}} (\emph{bailgāṛī}) is an ox cart: \emph{bail}, \textquotedblleft{}ox\textquotedblright{}, which CDIAL lists as Hindi \emph{bail} under \emph{*balilla}, \textquotedblleft{}ox\textquotedblright{} (CDIAL 9175), and \emph{gāṛī}, \textquotedblleft{}cart\textquotedblright{}, from \emph{*gāḍḍa} (CDIAL 4116).
\end{cousinweb}

\begin{grammarlens}[title={What you've built}]
One more word for this chapter.
\end{grammarlens}

\subsection*{Your turn}
\begin{itemize}
\raggedright
  \item \emph{Say it:} \emph{bailgāṛī}
  \item \emph{Say it:} \emph{bailgāṛī}, once more
  \item \emph{Say it:} say \emph{bailgāṛī}
  \item \emph{Recall:} say the Marwadi for an axe, then the Marwadi for a plough, then say \emph{bailgāṛī} again
\end{itemize}

\begin{tcolorbox}[breakable,colback=teal!4,colframe=teal!35!black,title={Before you move on}]
What does \mw{बैलगाड़ी} mean? (\textquotedblleft{}A bullock cart\textquotedblright{}.) Say it once more.
\end{tcolorbox}
